\documentclass[tc, manuscript]{copernicus}

\graphicspath{{../}}

\begin{document}

\title{A field-measurable regularized-Coulomb sliding law: the s\_N(N) master curve, an N\_c inversion, and a tidal-admittance ungrounding early-warning}

% Corresponding-author email: author's AIET institutional account.
\Author[1]{Abdelrahman}{Saifelden}

\affil[1]{Alexandria Higher Institute of Engineering and Technology (AIET), Alexandria, Egypt}

\runningtitle{A field-measurable regularized-Coulomb sliding law}
\runningauthor{A. Saifelden}
\correspondence{Abdelrahman Saifelden (Abdulrahman.Saifelden22011411@aiet.edu.eg)}

\received{}
\pubdiscuss{}
\revised{}
\accepted{}
\published{}

\firstpage{1}
\maketitle

\begin{abstract}
Ice-stream sliding weakens as the bed approaches flotation, a behaviour large-scale models impose by hand (Joughin, Smith \& Schoof 2019 add an ad-hoc near-flotation ramp to a regularized-Coulomb law). We show this weakening is a \textbf{closed-form pole} of the regularized-Coulomb law itself and, crucially, that it is \textbf{measurable from surface velocity alone}. Solving $\tau _{b}=C N (u/(u+u_{0}))^{1/m}$ at the driving stress gives, exactly, the sliding-sensitivity master curve $|s_{N}|(N)=m/(1-(N_{c}/N)^{m})$ with $N_{c}=\tau _{d}/C$ — a simple pole at the flotation threshold $N_{c}$, verified to $1.4\times 10^{-4}$ against the numeric $s_{N}$. The curve makes three independent surface observables read the \textbf{same} law: a lake-drainage step ($du/u=|s_{N}|\cdot dN/N$), a continuous ocean-thermal-forcing gating slope, and the \textbf{tidal admittance} of velocity (whose 2f/1f harmonic ratio diverges toward $N_{c}$). A drainage-step population inverts for $N_{c}$ to \textbf{0.2/0.4/1.0 \%} at 5/10/20 \% amplitude noise; the tidal probe recovers $m$ and the dimensionless flotation proximity $R=(N_{c}/N)^{m}$ from velocity alone (answering the ``no knowledge of $N$'' objection), with the tidal forcing amplitude \emph{measured}, not assumed (CATS2008 across the BedMachine grounding zone: median tidal pressure swing 0.30 \% of overburden, rising to 0.45 \% within 2–4 km of the line). Because the basal-drag stiffness $\partial \tau _{b}/\partial u\propto (1-R)^{2}/R\to 0$ at $N_{c}$, an ice stream nearing ungrounding should show rising variance and lag-1 autocorrelation of surface speed; we \textbf{operationalize} the established critical-slowing-down early-warning (Scheffer et al. 2009; Dakos et al. 2008; Rosier et al. 2021) as the continuous reading $R(t)=(N_{c}/N)^{m}\to 1$ from tidal admittance. Read through the same lens, three real MacAyeal-lake drainages sit far from the $N_{c}$ pole and show no critical-slowing-down signature — a correct \textbf{true negative} on stable ice (the detectors fire on planted surges/CSD). The contribution is the field-measurable law and its inversion, not a new early-warning concept; absolute $s_{N}$ remains uncalibrated and the strong test awaits a stream actually nearing flotation.
\end{abstract}

\introduction

The basal sliding law sets how fast an ice stream flows for a given driving stress and effective pressure $N=p_{i}-p_{w}$ (Schoof 2005; Gagliardini et al. 2007). Near the grounding line $N\to 0$ and sliding becomes Coulomb-limited; regularized-Coulomb laws (Joughin, Smith \& Schoof 2019; Tulaczyk et al. 2000; Gudmundsson 2011) capture this but introduce the near-flotation weakening through an \emph{ad hoc} ramp and a fixed velocity scale $u_{0}$. The difficulty is that $N$ is rarely known: it depends on the unobserved subglacial water pressure.

The tools we use, not claim, are the regularized-Coulomb friction law (Schoof 2005; Joughin et al. 2019), the nonlinear tidal (MSf) velocity response of ice streams (Gudmundsson 2011), and critical-slowing-down early-warning theory, \emph{including its ice-sheet application} (Scheffer et al. 2009; Dakos et al. 2008; Rosier et al. 2021; Boers \& Rypdal 2021).

\textbf{What is new, stated against the closest prior work.} Joughin, Smith \& Schoof (2019) \emph{impose} near-flotation weakening with a hand-tuned ramp; we \emph{derive} it as a closed-form pole $|s_{N}|=m/(1-(N_{c}/N)^{m})$ of the regularized-Coulomb law and turn $N_{c}$ into a \textbf{measured} quantity. Rosier et al. (2021) established critical-slowing-down indicators (lag-1 autocorrelation, variance, Kendall-\ensuremath{\tau}) for grounding-line flux ahead of marine-ice-sheet tipping; we do not re-claim the early-warning but give a specific \textbf{operational route} — reading $R(t)=(N_{c}/N)^{m}\to 1$ from the continuous tidal admittance of \emph{surface velocity}, which needs no knowledge of $N$.

\section{The $s_{N}(N)$ master curve}

\textbf{2.1 $s_{N}(N)$ as the common observable.} A lake drainage is an in-situ effective- pressure \emph{step} whose surge amplitude $du/u=|s_{N}|\cdot (dN/N)$ measures the sliding sensitivity $s_{N}=d ln u_{b}/d ln N$; ocean thermal forcing lowers $N$ \emph{continuously}, so the gating slope $d ln u_*/dTF$ measures the \textbf{same} $s_{N}$ ($+0.46\to +0.62/^{\circ}C$, direct BedMachine $N$); ocean tides oscillate $N$ at known amplitude, a third probe (§4). All three read one curve.

\textbf{2.2 The closed-form master curve.} Solving the regularized-Coulomb law $\tau _{b}=C N (u/(u+u_{0}))^{1/m}$ at $\tau _{b}=\tau _{d}$ gives, exactly, $|s_{N}|(N)=m/(1-(N_{c}/N)^{m})$, $N_{c}=\tau _{d}/C$ — $\to m$ far from flotation, a \textbf{simple pole} at $N_{c}$ ($\approx N_{c}/(N-N_{c})$). Verified to $1.4\times 10^{-4}$ against the numeric $s_{N}$. This \emph{derives} the near-flotation weakening that Joughin, Smith \& Schoof (2019) impose by hand (ad hoc $h_{af}<h_{T}$ ramp + fixed $u_{0}$). (\texttt{validation/\allowbreak{}synthetic/\allowbreak{}sn\_\allowbreak{}master\_\allowbreak{}curve.\allowbreak{}py})

\section{Inversion — measuring $N_{c}$ and $m$}

A drainage-step population recovers the flotation threshold $N_{c}$ to \textbf{0.2 / 0.4 / 1.0 \%} at 5 / 10 / 20 \% amplitude noise; $m$ is degenerate with the per-event drop far from flotation. The recovery is unbiased and sign-faithful under noise and null under permutation (synthetic plant-and-recover).

\section{The tidal third probe — $N_{c}$, $m$, and $R$ from velocity alone}

Ocean \textbf{tides} give a continuous probe: the fundamental velocity admittance equals $|s_{N}|$ and the 2f/1f harmonic ratio diverges toward $N_{c}$, so admittance + harmonics + the known tidal amplitude recover $m$ and the dimensionless flotation proximity $R=(N_{c}/N)^{m}$ \textbf{from surface velocity alone} — a sliding-law reading of the Gudmundsson nonlinear MSf response, answering the ``no knowledge of $N$'' objection. The tidal forcing amplitude is \textbf{measured}, not assumed: sampling the CATS2008 tide model (USAP-DC 601235) across the BedMachine grounding zone (within 10 km of the line; $n=1108$ valid cells) gives a tidal ocean-pressure swing with median \textbf{0.30 \% of ice overburden} (median tide $\eta \approx 0.84 m$), rising monotonically from 0.21 \% at 8–10 km to 0.45 \% at 2–4 km from the grounding line (\texttt{validation/\allowbreak{}synthetic/\allowbreak{}tidal\_\allowbreak{}admittance\_\allowbreak{}probe.\allowbreak{}py}, \texttt{validation/\allowbreak{}reports/\allowbreak{}tidal\_\allowbreak{}forcing\_\allowbreak{}gz.\allowbreak{}json}).

\section{An ungrounding early-warning (operationalizing established CSD theory)}

The $N_{c}$ fold makes the basal drag stiffness $\partial \tau _{b}/\partial u\propto (1-R)^{2}/R\to 0$ (critical slowing down), so an ice stream approaching ungrounding should show \textbf{rising variance and lag-1 autocorrelation of its surface speed} (Kendall-\ensuremath{\tau}\ensuremath{\approx}0.54 in an OU demonstration with $N$ declining toward $N_{c}$). Critical-slowing-down early-warning is established (Scheffer et al. 2009; Dakos et al. 2008) and has already been applied to marine-ice-sheet tipping — Rosier et al. (2021) detected rising lag-1 autocorrelation and variance in grounding-line flux ahead of Pine Island tipping, and Boers \& Rypdal (2021) in Greenland surface melt. The contribution here is \textbf{not} the early-warning concept but a specific \textbf{operational route}: reading $R(t)=(N_{c}/N)^{m}\to 1$ from the continuous tidal admittance of surface velocity. A single-snapshot \textbf{spatial} analogue (variance + along-flow correlation length rising toward the grounding line) is also derived (exact Lyapunov solve). (\texttt{validation/\allowbreak{}synthetic/\allowbreak{}spatial\_\allowbreak{}ews.\allowbreak{}py})


\textbf{The fold exponents and a calibration-free spectral gauge (derived).} The warning here is not generic: the $s_{N}$ pole \emph{is} the Schoof (2007) grounding-line saddle-node, so the critical-slowing-down exponents are fixed, not fitted. Near $N=N_{c}$ the restoring eigenvalue $\lambda \propto (1-R)^{2}/R\sim (N-N_{c})^{2}$, giving $|s_{N}|\propto (N-N_{c})^{-1}$, variance $\propto (N-N_{c})^{-2}$ and $AC1\to 1$ (recovered exponents $-1.00,\ -2.00,\ +2.00$). The same rate has a frequency-domain face: the velocity spectrum is a Lorentzian $S(\omega )=2D/(\lambda ^{2}+\omega ^{2})$ whose corner $f_{c}=\lambda /2\pi \propto (N-N_{c})^{2}$ \emph{reddens} toward ungrounding (the spectral early-warning of Kuehn 2011; Bury et al. 2020). Crucially the fluctuation--dissipation identity $\mathrm{Var}\cdot 2\pi f_{c}=D$ holds independent of $N$, so the noise corner $f_{c}$ is a \textbf{calibration-free} flotation-proximity gauge (it falls as $(N-N_{c})^{2}$ regardless of the unknown noise strength $D$); combined with the tidal-admittance \emph{response} (§4) it even returns $D$ itself --- a response/fluctuation pair from one velocity record. (Derived consistency relations, verified on OU realizations: corner recovered to $3\%$, FDT product constant to $5\times 10^{-17}$; \texttt{nr3\_\allowbreak{}misi\_\allowbreak{}fold\_\allowbreak{}ews}, \texttt{nr8\_\allowbreak{}spectral\_\allowbreak{}fdt\_\allowbreak{}ews}.)

\section{The framework on the same real lakes}

Reading open CryoSat-2 drainage + ITS\_LIVE velocity (3 MacAyeal lakes) through this lens (\texttt{validation/\allowbreak{}external/\allowbreak{}lake\_\allowbreak{}lag\_\allowbreak{}sn\_\allowbreak{}ews.\allowbreak{}py}, 5 tests): across the 3 resolved drainage events $|\Delta v/v|\le 2 \%$, mixed sign, near the \textasciitilde{}1 \% year-to-year noise floor — \textbf{0 positive >2\ensuremath{\sigma} surge detections} — so via the master curve these trunk lakes sit \textbf{far from the $N_{c}$ pole}; and \textbf{no} lake shows the joint critical-slowing-down signature (rising variance \emph{and} high AC1), so the early-warning correctly returns a \textbf{true negative} on stable ice (unit tests confirm the detectors \emph{do} fire on planted surges/CSD, so the null is real). Both probes agree: far from ungrounding.

\section{Estimator calibration, scope, and falsification}

\textbf{7.1 Calibration.} Every estimator has a synthetic plant-and-recover harness in \texttt{pytest}: the master-curve closed form ($1.4\times 10^{-4}$), the $N_{c}$ inversion (0.2/0.4/1.0 \%), the tidal admittance/harmonic recovery, and the CSD detectors (fire on planted surges/CSD, null on stable OU).

\textbf{7.2 Scope and falsification.} Absolute $s_{N}$ is uncalibrated (per-event \texttt{dN/\allowbreak{}N} is a lumped-storage lower bound); the robust claims are the sign/shape and the \emph{threshold} $N_{c}$. With n\ensuremath{\approx}8 annual points per lake the CSD test is weak; a strong test needs the dense quarterly/GPS series and a stream actually nearing flotation. Falsifiers: a co-located \texttt{dN}+GPS population whose $|s_{N}(N)|$ does not follow $m/(1-(N_{c}/N)^{m})$; a stream nearing flotation with no variance/AC1 rise; tidal admittance/harmonics not steepening toward the grounding line.

\section{Reproduce}

\begin{verbatim}
pip install -r requirements.txt
python glaciers/validation/synthetic/sn_master_curve.py        # §2 s_N(N) master curve + inversion
python glaciers/validation/synthetic/tidal_admittance_probe.py # §4 tidal third probe
python glaciers/validation/synthetic/spatial_ews.py            # §5 spatial early-warning
python glaciers/validation/external/lake_lag_sn_ews.py         # §6 framework on real lakes
pytest glaciers/tests/test_validation_synthetic.py -v          # §7 estimator calibration
\end{verbatim}

\codedataavailability{All analysis code and derived products are in the public repository at https://github.com/abdh0saifelden2-debug/K, and the results regenerate from the scripts above. The study uses the following open datasets:

\begin{itemize}
\item BedMachine Antarctica v4 (Morlighem, 2025, https://doi.org/10.5067/POJQI54A45HX; see Morlighem et al., 2020).
\item CATS2008 circum-Antarctic tide model: USAP-DC 601235 (Howard et al., 2019, https://doi.org/10.15784/601235).
\item ITS\_LIVE surface velocities (Gardner et al., 2025, https://doi.org/10.5067/JQ6337239C96; see Gardner et al., 2018) and CryoSat-2 lake-drainage dates.
\end{itemize}

BedMachine and ITS\_LIVE were obtained from the NASA NSIDC DAAC; CATS2008 from the U.S. Antarctic Program Data Center (https://www.usap-dc.org).}

\authorcontribution{A. Saifelden conceived the study, performed all computations and analysis, and wrote the manuscript.}

\competinginterests{The author declares that there are no competing interests.}

\begin{thebibliography}{}
\bibitem{ref1}
Boers, N., \& Rypdal, M. (2021). Critical slowing down suggests that the western Greenland Ice Sheet is close to a tipping point. \emph{Proceedings of the National Academy of Sciences} 118(21), e2024192118. doi:10.1073/pnas.2024192118

\bibitem{ref2}
Dakos, V., Scheffer, M., van Nes, E. H., Brovkin, V., Petoukhov, V., \& Held, H. (2008). Slowing down as an early warning signal for abrupt climate change. \emph{Proceedings of the National Academy of Sciences} 105(38), 14308–14312. doi:10.1073/pnas.0802430105

\bibitem{ref3}
Gagliardini, O., Cohen, D., Råback, P., \& Zwinger, T. (2007). Finite-element modeling of subglacial cavities and related friction law. \emph{Journal of Geophysical Research: Earth Surface} 112, F02027. doi:10.1029/2006JF000576

\bibitem{ref4}
Gardner, A. S., Moholdt, G., Scambos, T., Fahnestock, M., Ligtenberg, S., van den Broeke, M., \& Nilsson, J. (2018). Increased West Antarctic and unchanged East Antarctic ice discharge over the last 7 years. \emph{The Cryosphere} 12, 521–547. doi:10.5194/tc-12-521-2018

\bibitem{ref5}
Gardner, A., Fahnestock, M., Greene, C. A., Kennedy, J. H., Liukis, M., Lopez, L., \& Scambos, T. (2025). MEaSUREs ITS\_LIVE Regional Glacier and Ice Sheet Surface Velocities, Version 2 (NSIDC-0776) [Data set]. \emph{NASA NSIDC DAAC.} doi:10.5067/JQ6337239C96

\bibitem{ref6}
Gudmundsson, G. H. (2011). Ice-stream response to ocean tides and the form of the basal sliding law. \emph{The Cryosphere} 5, 259–270. doi:10.5194/tc-5-259-2011

\bibitem{ref7}
Howard, S. L., Erofeeva, S., \& Padman, L. (2019). CATS2008: Circum-Antarctic Tidal Simulation version 2008. \emph{U.S. Antarctic Program Data Center.} doi:10.15784/601235

\bibitem{ref8}
Joughin, I., Smith, B. E., \& Schoof, C. G. (2019). Regularized Coulomb friction laws for ice sheet sliding: Application to Pine Island Glacier, Antarctica. \emph{Geophysical Research Letters} 46, 4764–4771. doi:10.1029/2019GL082526

\bibitem{ref9}
Morlighem, M., Rignot, E., Binder, T., Blankenship, D. D., Drews, R., Eagles, G., et al. (2020). Deep glacial troughs and stabilizing ridges unveiled beneath the margins of the Antarctic ice sheet. \emph{Nature Geoscience} 13, 132–137. doi:10.1038/s41561-019-0510-8

\bibitem{ref10}
Morlighem, M. (2025). MEaSUREs BedMachine Antarctica, Version 4 (NSIDC-0756) [Data set]. \emph{NASA National Snow and Ice Data Center DAAC.} doi:10.5067/POJQI54A45HX

\bibitem{ref11}
Padman, L., Fricker, H. A., Coleman, R., Howard, S., \& Erofeeva, L. (2002). A new tide model for the Antarctic ice shelves and seas. \emph{Annals of Glaciology} 34, 247–254. doi:10.3189/172756402781817752

\bibitem{ref12}
Rosier, S. H. R., Reese, R., Donges, J. F., De Rydt, J., Gudmundsson, G. H., \& Winkelmann, R. (2021). The tipping points and early warning indicators for Pine Island Glacier, West Antarctica. \emph{The Cryosphere} 15, 1501–1516. doi:10.5194/tc-15-1501-2021

\bibitem{ref13}
Scheffer, M., Bascompte, J., Brock, W. A., Brovkin, V., Carpenter, S. R., Dakos, V., Held, H., van Nes, E. H., Rietkerk, M., \& Sugihara, G. (2009). Early-warning signals for critical transitions. \emph{Nature} 461, 53–59. doi:10.1038/nature08227

\bibitem{ref14}
Schoof, C. (2005). The effect of cavitation on glacier sliding. \emph{Proceedings of the Royal Society A} 461, 609–627. doi:10.1098/rspa.2004.1350

\bibitem{ref15}
Siegfried, M. R., \& Fricker, H. A. (2021). Illuminating active subglacial lake processes with ICESat-2 laser altimetry. \emph{Geophysical Research Letters} 48, e2020GL091089. doi:10.1029/2020GL091089

\bibitem{ref16}
Tulaczyk, S., Kamb, W. B., \& Engelhardt, H. F. (2000). Basal mechanics of Ice Stream B, West Antarctica: 1. Till mechanics. \emph{Journal of Geophysical Research: Solid Earth} 105(B1), 463–481. doi:10.1029/1999JB900329

\end{thebibliography}

\end{document}
